
\documentclass[11pt]{article}

\setlength{\parindent}{0pt}
\usepackage{xltxtra}
\usepackage{hyperref}
\hypersetup{hidelinks}
\usepackage{url}
\urlstyle{tt}
\usepackage{xcolor}
\definecolor{CVBlue}{RGB}{23,110,191}
\usepackage{calc}
\usepackage{graphicx}
\usepackage{tikz}
\usetikzlibrary{calc}
\usepackage{fontspec}
\usepackage{xeCJK}
\usepackage{enumitem}
\CJKsetecglue{}

\setmainfont[
    Path = fonts/Main/,
    Extension = .otf,
    BoldFont = texgyretermes-bold.otf
]{texgyretermes-regular.otf}

\setCJKmainfont[
    Path = fonts/hansans/,
    Extension = .ttf,
    BoldFont = NotoSansSC-Bold.ttf
]{NotoSansSC-Regular.otf}

\usepackage{relsize}
\usepackage{xspace}
\usepackage{fontawesome}

\usepackage[
a4paper,
left=1.2cm,
right=1.2cm,
top=1.5cm,
bottom=1cm,
nohead
]{geometry}

\renewcommand{\baselinestretch}{1.18}

\usepackage{titlesec}
\setlist{noitemsep}
\setlist[itemize]{topsep=0.15em,leftmargin=*}

\titleformat{\section}
{\large\bfseries\raggedright}
{}{0em}
{}
[{\color{CVBlue}\titlerule}]
\titlespacing*{\section}{0cm}{*1.25}{*0.8}

\begin{document}
\pagenumbering{gobble}

\begin{tikzpicture}[remember picture,overlay]
\node[anchor=north east] at ($(current page.north east)+(-2cm,-1.2cm)$)
{\includegraphics[height=3cm]{avatar.jpg}};
\end{tikzpicture}

\begin{tikzpicture}[remember picture,overlay]
\node[anchor=north west] at ($(current page.north west)+(0.5cm,+1.0cm)$)
{\includegraphics[height=5.5cm]{zju.png}};
\end{tikzpicture}

\centerline{\LARGE\bfseries 陈浩然}

\centerline{\normalsize{
\faPhone\ 15964026329\quad
\faEnvelopeO\ \href{mailto:Chenhr1221@zju.edu.cn}{Chenhr1221@zju.edu.cn}
\quad 籍贯：山东济南}}

\vspace{-0.2em}

\section{\makebox[\widthof{\faGraduationCap}][c]{\color{CVBlue}\faGraduationCap}\ 教育背景}

\textbf{浙江大学 工程师学院}
\hfill 2025.09 -- 至今

城市工程系统可持续项目 硕士研究生

\begin{itemize}
\item 研究方向：\textbf{工程结构安全、多灾种动力响应、数值模拟与工程数据分析}
\end{itemize}

\textbf{合肥工业大学（211/双一流） 土木与水利工程学院}
\hfill 2021.09 -- 2025.06

城市地下空间工程专业

\begin{itemize}
\item 专业综合排名：\textbf{11/249（前5\%）}，GPA：3.66/4.30，CET-6：457。
\item 主修课程：材料力学、弹性力学及有限元法、结构设计方法、概率论与数理统计。
\end{itemize}


\section{\makebox[\widthof{\faUsers}][c]{\color{CVBlue}\faUsers}\ 科研与项目经历}

\textbf{国家重点研发计划项目课题：强震和极端气候作用下灾害链动力学机制与对水电工程的影响预测}
\hfill 2025.09 -- 至今

项目成员

\begin{itemize}
\item 面向极端环境下工程灾害风险问题，开展\textbf{多灾种灾害链演化机制与动力响应分析}研究。
\item 基于\textbf{多源数据融合、Python数据处理与空间分析方法}，构建灾害演化数据库，开展工程风险因素分析。
\item 基于\textbf{3DEC离散元方法}开展三维数值模型构建，分析岩体结构响应及失稳演化过程。
\item 发明专利：\textbf{一种冰川灾害链物源识别与动态供源评价方法及系统}（第二发明人）。

\end{itemize}

\textbf{省级大学生创新创业训练计划：大学力学课程辅助学习系统开发}
\hfill 2023.05 -- 2024.05

核心成员

\begin{itemize}
\item 基于\textbf{MATLAB App Designer}开发力学课程辅助学习平台，实现知识点检索、典型例题管理及交互式学习功能。
\item 负责软件功能设计与界面开发，项目成果获得挑战杯校赛三等奖、互联网+校级铜奖，并完成软件著作权登记。
\end{itemize}


\section{\makebox[\widthof{\faCogs}][c]{\color{CVBlue}\faCogs}\ 工程实践经历}

\textbf{中交建筑集团第九工程有限公司}
\hfill 2024.06 -- 2024.09

项目技术岗实习生

\begin{itemize}
\item 参与房建项目技术管理工作，负责施工图纸整理、工程资料处理及项目技术支持。
\item 使用\textbf{AutoCAD、BIM}开展工程模型构建、图纸处理及数字化辅助分析。
\end{itemize}


\section{\makebox[\widthof{\faList}][c]{\color{CVBlue}\faList}\ 荣誉奖励}

\begin{itemize}
\item 全国大学生测绘学科创新创业智能大赛——无人机航测虚拟仿真 \textbf{一等奖}
\item 第十四届全国周培源大学生力学竞赛（安徽赛区）\textbf{二等奖}
\item 2023年高教社杯全国大学生数学建模竞赛（安徽赛区）\textbf{三等奖}
\item 第五届“设计总院杯”安徽省大学生桥梁设计大赛\textbf{一等奖}
\end{itemize}


\section{\makebox[\widthof{\faUsers}][c]{\color{CVBlue}\faUsers}\ 学生工作经历}

\begin{itemize}
\item 浙江大学研究生资助发展中心奖誉部干事（2025.09 -- 2026.08）；浙江大学工程师学院城市工程系统可持续项目生活委员、心理委员（2025.09 -- 至今）。
\end{itemize}


\section{\makebox[\widthof{\faCogs}][c]{\color{CVBlue}\faCogs}\ 技术技能}

\begin{itemize}
\item \textbf{CAE与数值模拟：}ANSYS有限元分析（FEA）、3DEC离散元模拟、MATLAB数值计算。
\item \textbf{工程设计与建模：}AutoCAD、BIM建模、工程数字化建模。
\item \textbf{编程与数据分析：}Python、SQL、Git、数据处理与可视化。
\item \textbf{GIS与遥感：}ArcGIS Pro、多源数据处理、空间统计分析。
\end{itemize}

\end{document}
